\documentclass[article]

\begin{document}


interpretation of evidence of LD instead of the few lines at the bottom of page 4. Maybe only as a part of the final discussion since it is not what the paper is about.

It is correct that we expect response dependence to be positive, but I do have examples where it is negative. E..g. because respondents interpret items as representing competing responses or because of the way that the questions are formulated. Evidence of negative local dependence also turns up because of simple item misfit. The point is that one should never interpret evidence of LD without looking at the items and at the tests of item misfit. In exactly the same way that one should never interpret evidence of item misfit without testing for LD. 

And, by the way, multidimensionality and trait dependence is supposed to generate evidence of both positive and negative local dependence. Positive dependence between items depending on the same latent trait and negative dependence between items associated with different traits. Stout’s exploratory DETECT procedure and the cDETECT procedure that Karl and I wrote about in our 2004 paper makes that point clear. It is for these reasons that I in most cases prefer two-sided tests of LD.

I am of course happy with the response dependence mechanism because it generates a nice log-linear Rasch model. However, I would also like to have a LD mechanism, which is a little more challenging to the gamma and the Q3. All the proposed LD mechanisms produce LD which is monotonic increasing. What about a mechanism with monotonic non-decreasing association. This actually often happens and is the reason why Kelderman’s test for LD appears to be much more powerful than the Gamma. Since Q3 appears to be a little more powerful than the Gamma, it would be interesting to see whether that also is true in such cases.  If it is not too late, I could try to define such a mechanism if it is not too late. I have simulated a phq9 data with your item parameters with n=600 (and n = 6000, just to see what happens) and could use that to define an appropriate mechanism for you.

Obviously, I expect it to be too late, so just say so. It is kind of fun playing with Rasch models in this way, so I will just see what happens with my simulated data.

Finally, I may have part of an explanation why Q3 may be more powerful than Gamma in small and moderate samples. The Q3 measure the association conditionally given scores from 0 to 27, whereas Gamma measure association conditionally given rest scores from 0-24. Gamma therefore excludes more persons with extreme scores than Q3. Whether that is the only reason, I have no idea. But it has to mean something with short scales and small samples.  In my simulated phq6000  data, Q3 disclosed spurious evidence of local dependence. That was expected because of the sample size, but Gamma could not see it.


\end{document}